\chapter{Criterio de publicación selectiva de claims prioritarios}

\section{Publicación selectiva}
\begin{definition}
Se llama publicación selectiva de claims prioritarios al procedimiento disciplinado mediante el cual un subconjunto de claims ya priorizados y suficientemente estabilizados se prepara para una forma externa de exposición más explícita, sobria y defendible.
\end{definition}

\begin{observation}
La publicación selectiva no equivale a apertura indiscriminada del corpus. Su función es declarar qué claims pueden sostener una primera exposición externa sin comprometer la sobriedad metodológica alcanzada en Tomo III.
\end{observation}

\section{Condición mínima de publicación selectiva}
\begin{definition}
Un claim prioritario puede entrar en publicación selectiva cuando concurren simultáneamente:
\begin{enumerate}
  \item prioridad curatorial ya justificada;
  \item compatibilidad con auditoría mínima;
  \item trazabilidad suficientemente limpia;
  \item alcance sobrio y legible;
  \item ausencia de reservas críticas no declaradas;
  \item capacidad de sostener una presentación externa sin extrapolación abusiva.
\end{enumerate}
\end{definition}

\begin{observation}
No se exige todavía una clausura total del sistema. Se exige solo que el claim pueda sostener una primera salida externa disciplinada sin deformar su fuerza metodológica real.
\end{observation}

\section{Publicación selectiva y tabla madre}
\begin{definition}
La publicación selectiva debe poder apoyarse en la tabla madre de demos y claims, de modo que toda afirmación publicada conserve vínculo visible con su origen, su prioridad y su estado metodológico previo.
\end{definition}

\begin{observation}
Esta cláusula evita que la publicación externa borre la genealogía interna del claim o lo haga parecer desligado de su proceso curatorial.
\end{observation}

\section{Publicación selectiva y anexo observado}
\begin{definition}
La publicación selectiva no puede ignorar la existencia de claims observados o reservados. Cuando un claim publicado tenga vínculos metodológicos relevantes con material observado, esa relación debe poder declararse con sobriedad.
\end{definition}

\begin{observation}
La salida externa no debe producir una imagen falsa de pureza total allí donde el propio proceso metodológico ha conservado reservas o tensiones abiertas.
\end{observation}

\section{Condición mínima de no publicación}
\begin{definition}
Un claim no debe entrar todavía en publicación selectiva cuando su prioridad sea solo interna, su fragilidad siga siendo alta o su forma comunicable externa exceda la fuerza metodológica realmente consolidada.
\end{definition}

\begin{observation}
La no publicación selectiva no implica descarte. Puede significar simplemente que el claim debe seguir en auditoría, reformulación o reserva observada.
\end{observation}

\section{Cláusula de alcance}
\begin{observation}
El criterio de publicación selectiva de claims prioritarios fija una primera regla de exposición externa sobria dentro del Tomo III. No constituye todavía una política editorial completa ni una apertura integral de todo el corpus.
\end{observation}

